% TOP_PROBLEM: {"problemId": "problem.top500-omission-w045049-potential-density-on-k3-surfaces", "canonicalTitle": "Potential Density of Rational Points on K3 Surfaces", "releaseRank": 483, "primaryDomain": "number_theory_arithmetic_geometry", "primaryDomainLabel": "Number theory & arithmetic geometry", "displayStatus": "open_with_solved_subcases", "releaseStatus": "open_with_solved_subcases"}
% !TeX program = lualatex
% Definition and sources only; this file is not an LLM solution attempt.
\documentclass[11pt]{article}
\usepackage[margin=25mm]{geometry}
\usepackage{fontspec}
\setmainfont{Latin Modern Roman}
\usepackage{amsmath,amssymb,mathtools}
\usepackage{unicode-math}
\setmathfont{Latin Modern Math}
\usepackage{xurl,hyperref}
\hypersetup{colorlinks=true,urlcolor=blue}
\setcounter{secnumdepth}{0}
\setlength{\parindent}{0pt}
\setlength{\parskip}{5pt}
\setlength{\emergencystretch}{3em}
\begin{document}
\section*{\texorpdfstring{Potential Density of Rational Points on K3 Surfaces}{Potential Density of Rational Points on K3 Surfaces}}\label{483}
\subsection{Definitions and mathematical statement}
Let $K$ be a number field. A K3 surface over $K$ is a smooth projective geometrically integral surface $X$ with trivial canonical line bundle and $H^1(X,\mathcal O_X)=0$. A subset is Zariski dense when it is not contained in any proper Zariski-closed subset. The conjecture is
\[
 \forall X/K\text{ K3}\quad\exists L/K\text{ finite}:\quad
 \overline{X(L)}^{\rm Zar}=X.
\]
The extension $L$ may depend on $X$ and $K$, with no uniform degree bound required. Density over the original field $K$ is not asserted. Infinitely many points on one algebraic curve inside $X$ need not be dense in the surface.
\par\smallskip\textit{Formulation and concept references:} \href{https://arxiv.org/html/2505.13262v2}{[S1]}.

\subsection{Short English statement}
Does every K3 surface over a number field acquire a Zariski-dense set of rational points after some finite field extension?
\subsection{Sources}
\begin{itemize}
\item[S1]J. Martínez-Marín, Rational points on K3 surfaces of degree 2, version 2 (2025), Introduction.
\url{https://arxiv.org/html/2505.13262v2}.
\textit{Formulation source.}
\end{itemize}
\end{document}
